\chapter{Inline images and graphics protocols}
\label{ch:images}

\begin{aside}
The terminal was always character-cell. Then someone wanted to
embed a thumbnail in their file browser. Two protocols compete
to make that work, plus one venerable hack from the 1970s.
\end{aside}

\section{Three protocols}

\begin{itemize}[leftmargin=*]
  \item \textbf{Sixel.} DEC's 1980s format: each pixel encoded
    as a six-bit byte representing a 1x6 column of pixels.
    Supported by xterm (compile flag), foot, mlterm, Wezterm.
  \item \textbf{iTerm2.} Image-as-base64 plus a custom escape.
    Supported by iTerm2 and Wezterm.
  \item \textbf{Kitty.} A chunked, multi-frame, palette-aware
    protocol. Supported by Kitty and Wezterm.
\end{itemize}

\maya{} can use any of the three; the detection sequence picks
the most capable.

\section{Sixel basics}

A sixel byte encodes six vertically stacked pixels. Bit 0 is
the top, bit 5 is the bottom. The character is \code{(byte + 63)}
to put it in the printable ASCII range.

A sixel string looks like:

\begin{lstlisting}
ESC P 7 ; 1 q "1;1;100;60 #0;2;100;100;100 #0!100? -- $
\end{lstlisting}

Where:
\begin{itemize}[leftmargin=*]
  \item \code{ESC P 7 ; 1 q}: enter sixel mode.
  \item \code{"1;1;100;60}: aspect and size.
  \item \code{\#0;2;100;100;100}: define colour 0 as RGB.
  \item Glyphs are pixel data; \code{!100?} is run-length
    encoding (100 copies of \code{?}).
  \item \code{-} starts a new sixel row; \code{\$} carriage
    return within a row.
\end{itemize}

It is grim to write by hand. \maya{} ships an encoder.

\section{iTerm2}

The iTerm2 protocol is simpler. Encode the image as base64 and
wrap:

\begin{lstlisting}
ESC ] 1337 ; File = inline=1 ; size=BYTES ; name=BASE64NAME :
       BASE64DATA BEL
\end{lstlisting}

Decode: the terminal renders the image at the current cursor
position, scaling to fit cell boundaries or to a specified
geometry.

This is easier to emit but offers less control: no animation,
no palette, no chunking.

\section{Kitty protocol}

The Kitty protocol allows everything: chunked transfers (split
large images across multiple escape sequences), placement IDs
(modify a placed image later), animations, palette images,
exact cell-bound placement.

A simple Kitty placement:

\begin{lstlisting}
ESC _ G a=T,f=32,s=WIDTH,v=HEIGHT,c=COLS,r=ROWS ;
       BASE64DATA ESC \
\end{lstlisting}

Where \code{a=T} means ``transmit and display'', \code{f=32}
is the format (RGBA), and \code{s,v} are the pixel
dimensions.

\maya{}'s image element wraps this:

\begin{lstlisting}
auto thumb = image(load("cat.png"))
           | width(20) | height(10);
\end{lstlisting}

\section{Cell-pixel sizing}

To place an image at a cell-grid boundary, you need to know
the terminal's cell size in pixels. \maya{} queries with
\code{ESC [ 16 t} (cell size) at startup; the reply is
\code{ESC [ 6 ; H ; W t}.

Without a reply, \maya{} falls back to a heuristic: 8x16 for
most fonts. Inline images may look off-grid; the user can
resize or set an explicit cell size in config.

\section{Z-stacking}

Inline images sit on top of cells. When the diff redraws a
cell that an image overlaps, the image must be re-emitted or
it will be clipped.

\maya{}'s renderer tracks image placements and re-emits them
after any frame that touches their region. For a static image
in a quiet UI, this is one emission. For an image surrounded
by changing cells, it is more.

\begin{warning}
Inline images are expensive. A 320x240 image is roughly 80 KB
of base64 per emission. Frequent re-emission is wire-heavy.
Place images in stable regions (e.g. fixed-position thumbnail
panels), not in scrolling areas.
\end{warning}

\section{Animations and APNG}

The Kitty protocol supports animations with frame chunks and
timing metadata. \maya{} can emit a multi-frame placement and
the terminal handles the animation.

The application emits frames; the terminal owns timing. This is
the opposite of how widget animations work (where \maya{} owns
timing). The reason: terminal-side animations don't wake the
application thread, so a spinner does not cost CPU.

\section{The unicode block trick}

For terminals without any image protocol, a low-res
approximation uses the half-block character \code{U+2580}:

\begin{itemize}[leftmargin=*]
  \item Foreground colour: the top half of the cell.
  \item Background colour: the bottom half.
  \item Glyph: \code{\u{2580}} (upper half block).
\end{itemize}

Each cell encodes two pixels. A 40x24-cell image is 40x48 pixels.

\maya{}'s \code{block\_image} element uses this. It looks
chunky but works on any terminal.

\begin{funfact}
The half-block trick is older than the web; it appeared in
TRS-80 and Apple II graphics modes via the same characters.
The terminal world reused the technique with the addition of
RGB colour.
\end{funfact}

\section{Image element API}

\begin{lstlisting}
auto img = image(path_or_buffer)
         | width(40) | height(12)    // in cells
         | fit(Fit::Contain);        // or Cover, Stretch
\end{lstlisting}

The fit modes:
\begin{itemize}[leftmargin=*]
  \item \textbf{Contain.} Scale to fit, preserve aspect; may
    leave empty cells.
  \item \textbf{Cover.} Scale to fill, preserve aspect; may
    crop.
  \item \textbf{Stretch.} Distort to fit exactly.
\end{itemize}

\section{Loading}

\maya{} ships PNG and JPEG decoders. SVG is not built in;
rasterise to PNG first with a tool. Animated GIF is decoded
to frames and looped via the protocol.

\section{Recap}

\begin{recap}
\begin{itemize}[leftmargin=*]
  \item Sixel, iTerm2, Kitty: three live protocols.
  \item Cell-pixel sizing via \code{ESC [ 16 t}.
  \item Re-emit images after diff overlaps.
  \item Half-block trick for terminals without any protocol.
  \item Fit modes: Contain, Cover, Stretch.
\end{itemize}
\end{recap}

\section*{Workshop}

\begin{enumerate}[leftmargin=*]
  \item Detect which image protocol your terminal supports.
    Display the result in your status bar.
  \item Build a thumbnail panel for an image directory.
  \item Toggle between Kitty protocol and the half-block
    fallback at runtime. Compare visual quality and wire
    cost.
\end{enumerate}
